\scanpage{333}
\mspage{4}

On vérifie aisément que

\[
 H^0(T,\mathscr H^q)=0\qquad\text{si }q\geq1
\]

et il reste

\[
 \begin{aligned}
 H^n(\deleted{\unsure{T}}X,M)&\simeq H^1(S,\mathscr H^{n-1})=\operatorname{Coker}\bigl(\delta:H^{n-1}(G)\longrightarrow{}\\
 &\hspace{12em}H^{n-1}(G\times G)\bigr)
 \end{aligned}
\]

Cela donne, pour $n=2$

\[
 \begin{aligned}
 H^2(X,M)&\simeq\operatorname{Coker}\bigl(\delta:\operatorname{Hom}_{\mathrm{mon}}(G,M)\longrightarrow
 \underbrace{\operatorname{Hom}_{\mathrm{mon}}(G\times G,M)}_{\operatorname{Hom}_{\mathrm{mon}}(G,M)\times\operatorname{Hom}_{\mathrm{mon}}(\unsure{G,M})}\bigr)\\
 &\simeq\operatorname{Coker}\bigl(u\longmapsto(u,-u):H\longrightarrow H\times H\bigr)
 \end{aligned}
\]

où $H$ est le groupe abélien $\operatorname{Hom}_{\mathrm{mon}}(G,M)$.

D'où

\[
 H^2(X,M)\simeq H=\operatorname{Hom}_{\mathrm{mon}}(G,M).
\]

Donc $X$ est 2-connexe \unsure{sss|ssi} lé groupe \uline{abélien} enveloppant du monoïde $G$ est nul.

Ainsi, si $G$ est un groupe abélien $\ne0$, $X$ n'\uline{est pas} 2-connexe.

\deleted{Conjecture} \uline{\unsure{Corollaire|Conclusion}} $\deleted{\unsure{\mathscr D}}\unsure{\Psi|\mathscr V|V}={}$

\diagram{\begin{tikzcd}[ampersand replacement=\&,row sep=tiny]
 \& b_0\\
 a_0\arrow[ur]\arrow[dr] \&\\
 \& c_0
\end{tikzcd}}

n'est \uline{pas} \deleted{\unsure{tot.}} $W_\infty$-asphériseur ; i.e. si on a un diagramme

\diagram{\begin{tikzcd}[ampersand replacement=\&,row sep=tiny]
 \& b\\
 a\arrow[ur]\arrow[dr] \&\\
 \& c
\end{tikzcd}}

dans une cat. $X\in\mathrm{Ob}\,\mathrm{Cat}$, telle que \unsure{sa|de} somme amalgamée dans $X^\wedge$ soit égale à $e_X$, il ne s'ensuit pas que $X$ soit $W_\infty$-asphérique, (Mais \deleted{\unsure{semble}} peut-être $X$ \unsure{voudra|doive} bien être 1-connexe ?)

% archive sheet number: 333
